% !TeX program = xelatex
\documentclass[11pt,a4paper]{article}
\usepackage[utf8]{inputenc}
\usepackage[T1]{fontenc}
\usepackage[english]{babel}
\usepackage{amsmath,amssymb,amsfonts,mathtools,amsthm}
\usepackage{geometry}
\geometry{left=1.25cm,right=1.25cm,top=1.25cm,bottom=3.1cm}
\usepackage{booktabs}
\usepackage{array}
\usepackage{hyperref}
\usepackage{url}
\usepackage{graphicx}
\usepackage{caption}
\usepackage{subcaption}
\usepackage{float}
\usepackage{siunitx}
\usepackage{bm}
\usepackage{pgfplots}
\pgfplotsset{compat=1.18}
\usepackage{xcolor}
\usepackage{titlesec}
\usepackage{fancyhdr}
\usepackage{microtype}
\usepackage{fontspec}
\usepackage{tcolorbox}
\usepackage{enumitem}
\usepackage{xeCJK}  % 한글 지원

% 라틴 폰트 (원본과 동일)
\setmainfont{Calibri}
\setsansfont{Segoe UI}[
BoldFont = Segoe UI Bold,
Ligatures = TeX
]

% 한국어 폰트 (Windows 기본)
\setCJKmainfont{Malgun Gothic}
\setCJKsansfont{Malgun Gothic}
\setCJKmonofont{Malgun Gothic}

\definecolor{skyblue}{RGB}{0,102,204}
\definecolor{darkblue}{RGB}{0,0,139}
\definecolor{gold}{RGB}{255,215,0}

% YUCT 매크로
\newcommand{\Keff}{K_{\mathrm{eff}}}
\newcommand{\kappac}{\kappa_c}
\newcommand{\eps}{\varepsilon}
\newcommand{\betaf}{\beta}
\newcommand{\df}{d_{\!f}}
\newcommand{\Sodd}{S_{\mathrm{odd}}}
\newcommand{\Seven}{S_{\mathrm{even}}}
\newcommand{\Mpl}{M_{\mathrm{Pl}}}
\newcommand{\Lpl}{l_{\mathrm{Pl}}}
\newcommand{\tpl}{t_{\mathrm{Pl}}}
\newcommand{\Lmin}{\ell_{\min}}
\newcommand{\LUV}{\Lambda_{\mathrm{UV}}}

% 정리 환경 (한국어)
\newtheorem{theorem}{정리}[section]
\newtheorem{proposition}[theorem]{명제}
\newtheorem{definition}[theorem]{정의}
\theoremstyle{remark}
\newtheorem{remark}[theorem]{비고}

% 섹션 서식
\titleformat{\section}{\LARGE\bfseries\color{darkblue}}{\thesection}{1em}{}
\titleformat{\subsection}{\Large\bfseries\color{darkblue}}{\thesubsection}{1em}{}
\titleformat{\subsubsection}{\normalsize\bfseries\color{darkblue}}{\thesubsubsection}{1em}{}

% 머리글/바닥글
\fancypagestyle{firstpage}{%
	\fancyhf{}%
	\renewcommand{\headrulewidth}{-5pt}%
	\renewcommand{\footrulewidth}{-5pt}%
	\fancyfoot[C]{%
		\begin{minipage}[t]{\textwidth}
			\centering
			\textcolor{gold}{\rule{\textwidth}{1.6pt}}\\[5pt]
			{\small \textsuperscript{1}Yakushev Research, \url{https://yuct.org/}} \hfill
			{\small YPSDC 프로토콜: \url{https://ypsdc.com/}}\\[-2pt]
			\textcolor{gold}{\rule{\textwidth}{1.6pt}}\\[5pt]
			\textbf{야쿠셰프 통일 조정 이론 2026} \hfill \thepage\\[10pt]
		\end{minipage}%
	}%
}

\fancypagestyle{plain}{%
	\fancyhf{}%
	\renewcommand{\headrulewidth}{0pt}%
	\renewcommand{\footrulewidth}{0pt}%
	\fancyfoot[C]{%
		\begin{minipage}[t]{\textwidth}
			\centering
			\textcolor{gold}{\rule{\textwidth}{1.6pt}}\\[4pt]
			\textbf{야쿠셰프 통일 조정 이론 2026} \hfill \thepage\\[4pt]
		\end{minipage}%
	}%
}

\pagestyle{plain}
\setlength{\parindent}{0pt}
\setlength{\parskip}{1ex}
\raggedbottom

\hypersetup{
	colorlinks=true,
	linkcolor=blue,
	citecolor=blue,
	urlcolor=blue,
}

% 사용자 정의 헤더
\newcommand{\makeheader}{%
	\noindent
	\begin{minipage}{0.17\textwidth}
		\raggedright
		{\fontspec{Segoe UI}[BoldFont=Segoe UI Bold]\bfseries\color{skyblue}\fontsize{46}{46}\selectfont YUCT}%
	\end{minipage}%
	\hspace{30pt}
	\begin{minipage}{0.32\textwidth}
		\raggedright
		{\sffamily\fontsize{11}{13}\selectfont YAKUSHEV\\ UNIFIED\\ COORDINATION THEORY}%
	\end{minipage}%
	\hfill
	\begin{minipage}{0.38\textwidth}
		\raggedleft
		{\sffamily\bfseries 부록 UVD}\\ 
		{\sffamily 버전 1.0 / 2026년 6월}\\ 
		{\sffamily\footnotesize \url{https://doi.org/10.5281/zenodo.18444598}}%
	\end{minipage}
	\par\vspace{5pt}
	\noindent\textcolor{gold}{\rule{\textwidth}{1.6pt}}
	\par\vspace{0pt}
}

\begin{document}
	\thispagestyle{firstpage}
	\makeheader
	
	\vspace{0.5cm}
	{\LARGE\bfseries 부록 UVD: YUCT에서 자외선 발산의 해결 \\
		\Large 조정 형상 인자와 조정 네트워크의 물리적 규모로부터 얻는 유한 양자장론 \par}
	\vspace{0.3cm}
	
	{\large Alexey V. Yakushev\textsuperscript{1}} \\
	\vspace{0.2cm}
	
	{\color{darkblue}\textbf{\LARGE 초록}} \par
	{\large
		\noindent
		야쿠셰프 통일 조정 이론(YUCT)은 진공의 유한한 정보 용량에서 비롯된 조정 형상 인자 \(F(q^2) = \exp[-(q^2/\LUV^2)^{2/3}]\)에 의해 양자장론의 자외선 발산이 제거된다고 제안한다. 원래 정식화에서 자외선 규모 \(\LUV\)는 플랑크 질량 \(\LUV = 3\pi\Mpl \approx 1.15\times 10^{20}\) GeV로 동일시되어 힉스 질량에 대한 거대한 유한 복사 보정을 초래했다. 우리는 스칼라 질량 보정에 나타나는 상수 \(C_0\)의 계산 오류를 수정한다. 올바른 값은 \(C_0 = 3\sqrt{\pi}/(8\sqrt{2}) \approx 0.46999\)이며, 이전에 보고된 \(C_0 \approx 0.626\)을 대체한다. 이 수정으로 플랑크 규모 네트워크로부터 힉스 질량에 대한 복사 기여는 \(\delta m \approx 1.57\times 10^{18}\) GeV가 되어, 맨 질량이 정밀하게 상쇄되지 않으면 심각한 미세 조정 문제를 다시 도입한다. 우리는 두 가지 논리적 대안을 검토한다: (i) YUCT 질량 공식 \(m = \Mpl(2/3)^L\xi_f\)가 이미 물리적 질량을 결정하며, 복사 보정은 유한한 상쇄 항에 의해 흡수된다 – 기술적으로 일관되지만 미학적으로 만족스럽지 못한 해결책; (ii) 플랑크 규모의 기본 라그랑지안은 맨 질량 항을 포함하지 않으며(규모 불변성), 전체 힉스 질량이 복사 보정에 의해 생성된다. 이 두 번째 시나리오에서는 자기 일관성이 자외선 규모가 \(\LUV \approx 8.96\) TeV일 것을 요구하며, 이는 LHC에서 접근 가능한 값이다. 그러면 조정 형상 인자는 TeV 규모에서 고에너지 산란 단면적의 관측 가능한 편차를 예측한다. 우리는 두 시나리오와 그 수학적 정당성, 실험적 결과를 제시하여 YUCT를 엄격한 경험적 검증에 노출시키는 것을 목표로 한다.
	}
	
	\vspace{0.2cm}
	\noindent
	{\color{darkblue}\textbf{키워드:}} YUCT, 자외선 발산, 조정 형상 인자, 계층 문제, 힉스 질량, 규모 불변성, 비국소 양자장론, LHC 현상론.
	\vspace{0.1cm}
	\noindent
	{\color{darkblue}\textbf{인용:}} Yakushev AV. 부록 UVD: YUCT에서 자외선 발산의 해결. 2026. DOI: \url{https://doi.org/10.5281/zenodo.18444598} (이 권).
	
	\vspace{0.5cm}
	
	% ============================================================
	\section{서론}
	\label{sec:intro}
	
	야쿠셰프 통일 조정 이론(YUCT)은 양자장론(QFT)의 자외선 발산에 대한 우아한 해결책을 제공한다. 임의의 운동량 차단 또는 무한 재규격화를 도입하는 대신, YUCT는 물리적 진공이 양자 정보를 전송하기 위한 유한한 "대역폭"을 가진다고 가정한다. 이 대역폭은 특성 규모 \(\LUV\) 이상의 가상 운동량을 지수적으로 억제하는 \textbf{조정 형상 인자} \(F(q^2)\)에 의해 정량화된다. 이 형상 인자의 수학적 구조는 YUCT의 보편적 오차 법칙 \(\varepsilon \propto K_{\mathrm{eff}}^{-\beta}\) (\(\beta = 2/3\))에 의해 결정된다.
	
	이 부록의 이전 버전에서 우리는 \(\LUV\)를 플랑크 질량 \(\LUV = 3\pi\Mpl \approx 1.15\times 10^{20}\) GeV로 동일시했다. 이는 YUCT의 1차 대수적 루프에 의해 결정된 최소 길이 \(\Lmin = \frac{2}{3}\Lpl\)에서 유도된 규모이다. 이 동일시에 의해 모든 루프 적분은 명백히 유한해진다: 로그 발산은 \(\ln(\LUV/m)\)로 대체되고, 멱급수 발산은 완전히 사라진다.
	
	그러나 스칼라 질량 보정의 신중한 재계산은 유한 잔여물의 크기를 지배하는 상수 \(C_0\)의 계산 오류를 드러낸다. \(C_0\)의 올바른 값은 이전에 보고된 것보다 작아 복사 보정을 줄이지만, 여전히 관측된 힉스 질량보다 훨씬 큰 자릿수로 남는다. 이는 물리적 규모 \(\LUV\)에 대한 비판적 재검토를 강제한다.
	
	이 부록은 두 가지 목표를 가진다. 첫째, 우리는 수정된 수학적 유도를 완전한 세부 사항과 함께 제시하여 모든 표현이 정확하고 재현 가능하도록 한다. 둘째, 우리는 수정된 공식의 물리적 함의를 탐구한다. 여기에는 조정 네트워크의 기본 규모가 플랑크 규모가 아니라 훨씬 낮은 규모일 수 있으며, 현재 충돌 실험에서 도달 가능할 수 있다는 급진적 제안이 포함된다.
	
	% ============================================================
	\section{조정 형상 인자와 자외선 규모}
	\label{sec:formfactor}
	
	YUCT 조정 형상 인자는 다음과 같다:
	\begin{equation}
		F(q^2) = \exp\!\Bigl[-\Bigl(\frac{q^2}{\LUV^2}\Bigr)^{\!\beta}\,\Bigr], \qquad \beta = \frac{2}{3}.
		\label{eq:F}
	\end{equation}
	모든 표준 모형의 전파 인자는 \(F(p^2/\LUV^2)\)로 곱해진다. 규모 \(\LUV\)는 원래 콤프턴 관계 \(\LUV = 2\pi\hbar c/\Lmin = 3\pi\Mpl\)를 통해 최소 길이 \(\Lmin = \frac{2}{3}\Lpl\)에서 유도된다.
	
	이 형상 인자는 모든 루프 적분을 자외선 영역에서 절대 수렴하게 만든다. 증명은 간단하다: 루프 운동량의 임의의 다항식 \(P(k)\)에 대해, 곱 \(P(k) \cdot \prod_i F(k_i^2/\LUV^2)\)은 지수 함수가 지배적이므로 \(|k| \to \infty\)일 때 \(|k|\)의 어떤 역멱수보다 빠르게 감소한다.
	
	% ============================================================
	\section{전자 자기 에너지: 로그 발산의 해결}
	\label{sec:electron}
	
	QED에서 1-루프 전자 자기 에너지는 기존 QFT에서 로그 발산한다. YUCT 형상 인자를 사용하면 유클리드 적분은 다음과 같아진다:
	\begin{equation}
		\Sigma_{\mathrm{YUCT}} \propto \int_0^{\infty} \frac{dk_E}{k_E}\; \exp\!\Bigl[-2\Bigl(\frac{k_E^2}{\LUV^2}\Bigr)^{\!2/3}\,\Bigr].
		\label{eq:elec_int}
	\end{equation}
	치환 \(t = (k_E^2/\LUV^2)^{2/3}\)은 \(dk_E/k_E = \frac{3}{4} dt/t\)를 주고, 적분은 다음과 같이 평가된다:
	\begin{equation}
		I_{\mathrm{log}} = \frac{3}{4} \int_{t_{\min}}^{\infty} \frac{dt}{t}\, e^{-2t} \approx \ln\frac{\LUV}{m_e},
		\label{eq:elec_result}
	\end{equation}
	여기서 \(t_{\min} \sim (m_e^2/\LUV^2)^{2/3}\)이다. 로그 발산은 유한 로그 \(\ln(\LUV/m_e) \approx 53.8\)로 대체된다. 이 계산은 초등적이며 수정이 필요하지 않다.
	
	% ============================================================
	\section{스칼라 질량 보정: 상수 \(C_0\)의 수정}
	\label{sec:scalar}
	
	\subsection{적분 및 그 평가}
	
	\(\phi^4\) 이론에서 1-루프 스칼라 질량 보정은, 위크 회전과 형상 인자 삽입 후,
	\begin{equation}
		\delta m^2_{\mathrm{YUCT}} = \frac{\lambda}{16\pi^2} \int_0^{\infty} \frac{k_E^3\, dk_E}{k_E^2 + m^2}\; F^2\!\Bigl(\frac{k_E^2}{\LUV^2}\Bigr), \qquad F^2(x) = e^{-2x^{2/3}}.
		\label{eq:scalar_YUCT}
	\end{equation}
	이다.
	
	무차원 변수 \(t = k_E^2/\LUV^2\)를 도입하고 \(m \ll \LUV\)에 대해 전개하면, 주도 항은
	\begin{equation}
		\delta m^2_{\mathrm{YUCT}} = \frac{\lambda}{32\pi^2}\,\LUV^2 \int_0^{\infty} e^{-2t^{2/3}} dt + \mathcal{O}(m^2/\LUV^2).
		\label{eq:scalar_lead}
	\end{equation}
	이다.
	
	이전(버전 3.1)에 적분 \(C_0 \equiv \int_0^{\infty} e^{-2t^{2/3}} dt \approx 0.626\)으로 보고되었다. 이 값은 틀렸다. 이제 정확히 계산한다.
	
	\subsection{\(C_0\)의 정확한 계산}
	
	\begin{theorem}[\(C_0\)의 정확한 값]
		\label{thm:C0}
		\[
		C_0 = \int_0^{\infty} e^{-2t^{2/3}} dt = \frac{3\sqrt{\pi}}{8\sqrt{2}} \approx 0.46999.
		\]
	\end{theorem}
	
	\begin{proof}
		\(u = 2t^{2/3}\)로 치환하자. 그러면 \(t = (u/2)^{3/2}\)이고,
		\[
		dt = \frac{3}{2} (u/2)^{1/2} \cdot \frac{1}{2} du = \frac{3}{4\sqrt{2}} u^{1/2} du.
		\]
		따라서
		\[
		C_0 = \int_0^{\infty} e^{-u} \cdot \frac{3}{4\sqrt{2}} u^{1/2} du = \frac{3}{4\sqrt{2}} \int_0^{\infty} u^{1/2} e^{-u} du = \frac{3}{4\sqrt{2}} \, \Gamma(3/2).
		\]
		\(\Gamma(3/2) = \frac{\sqrt{\pi}}{2}\)을 사용하면,
		\[
		C_0 = \frac{3}{4\sqrt{2}} \cdot \frac{\sqrt{\pi}}{2} = \frac{3\sqrt{\pi}}{8\sqrt{2}} \approx 0.46999.
		\]
	\end{proof}
	
	\subsection{수정된 스칼라 질량 보정}
	
	올바른 \(C_0\) 값을 사용하면, 스칼라 질량 보정의 주도 항은
	\begin{equation}
		\boxed{\delta m^2_{\mathrm{YUCT}} = \frac{\lambda}{32\pi^2}\,\LUV^2 \cdot \frac{3\sqrt{\pi}}{8\sqrt{2}} + \mathcal{O}(m^2/\LUV^2)}.
		\label{eq:scalar_corrected}
	\end{equation}
	이다.
	
	힉스 보손에 대해 \(\lambda = m_H^2/(2v^2) \approx 0.13\) (여기서 \(v \approx 246\) GeV)이다. 플랑크 규모 동일시 \(\LUV = 3\pi\Mpl \approx 1.15\times 10^{20}\) GeV를 사용하면:
	\begin{align}
		\delta m^2_H &\approx \frac{0.13}{32\pi^2} \cdot (1.15\times 10^{20})^2 \cdot 0.46999 \notag\\
		&\approx 2.47 \times 10^{36}\ \mathrm{GeV}^2, \notag\\
		\delta m_H &\approx 1.57 \times 10^{18}\ \mathrm{GeV}.
		\label{eq:planck_correction}
	\end{align}
	
	이 값은 관측된 힉스 질량 (\(m_H^{\mathrm{exp}} \approx 125\) GeV)보다 \(1.26 \times 10^{16}\)배 더 크다. 라그랑지안의 맨 질량이 0이라면 물리적 질량은 전적으로 이 복사 보정에 의해 결정되며, 예측된 값은 명백히 관측과 양립할 수 없다.
	
	% ============================================================
	\section{물리적 함의: 두 시나리오}
	\label{sec:scenarios}
	
	수정된 계산은 계층 문제에 직접 직면하도록 강제한다. 우리는 논리적으로 일관된 두 가지 대안을 검토한다.
	
	\subsection{시나리오 A: 유한 상쇄 항을 갖는 플랑크 규모 네트워크}
	
	이 시나리오에서 YUCT는 동일시 \(\LUV = 3\pi\Mpl\)를 유지한다. 임의 입자의 물리적 질량은 YUCT 질량 공식
	\begin{equation}
		m_{\mathrm{phys}} = \Mpl \left(\frac{2}{3}\right)^L \xi_f,
		\label{eq:mass_formula}
	\end{equation}
	에 의해 주어진다. 여기서 \(L\)은 조정 깊이, \(\xi_f\)는 공명 인자이다(부록~J 참조). 라그랑지안의 맨 질량은 자유 매개변수로, 복사 보정을 상쇄하고 관측된 질량을 재현하도록 선택된다.
	
	수학적으로 이것은 완벽하게 일관된다: 이론은 유한하며, 필요한 상쇄 항도 유한하고, 무한대는 결코 발생하지 않는다. YUCT 질량 공식 Eq.~(\ref{eq:mass_formula})는 상쇄 항을 고정하는 재규격화 조건을 제공한다.
	
	그러나 물리적으로 이 시나리오는 맨 질량과 복사 보정이 \(10^{16}\) 분의 1의 정밀도로 상쇄될 것을 요구한다 – 이는 정확히 계층 문제가 피하려고 했던 종류의 미세 조정이다. YUCT는 왜 이 상쇄를 달성하기 위해 맨 질량을 선택해야 하는지 설명하지 않는다; 그것은 단지 일반적인 무한 뺄셈을 거대한 유한 뺄셈으로 대체할 뿐이다.
	
	\subsection{시나리오 B: 규모 불변 기원과 TeV 규모 네트워크}
	
	미적으로 더 매력적인 대안은 \(\LUV\)의 플랑크 규모 동일시를 포기하는 것이다. 만약 플랑크 규모의 기본 이론이 맨 질량 항을 갖지 않는다면(즉, 규모 불변이라면), 모든 질량은 조정 형상 인자에 의해 방사적으로 생성된다. 이 경우 물리적 힉스 질량 \emph{은} 복사 보정이다:
	\begin{equation}
		m_H^2 = \delta m^2_{\mathrm{YUCT}} = \frac{\lambda}{32\pi^2}\,\LUV^2 \cdot C_0.
		\label{eq:scale_inv}
	\end{equation}
	
	\(\LUV\)에 대해 풀면:
	\begin{equation}
		\LUV = \sqrt{\frac{32\pi^2\, m_H^2}{\lambda\, C_0}}.
		\label{eq:LUV_solve}
	\end{equation}
	
	수정된 \(C_0 = 3\sqrt{\pi}/(8\sqrt{2}) \approx 0.46999\), \(\lambda = 0.13\), 그리고 \(m_H = 125\) GeV를 대입하면,
	\begin{equation}
		\boxed{\LUV \approx 8.96\ \mathrm{TeV}}.
		\label{eq:LUV_TeV}
	\end{equation}
	를 얻는다.
	
	이 값은 이론의 예측이며 자유 매개변수가 아니다. 이것은 기본 조정 규모를 TeV 범위에 정확히 위치시켜 실험적으로 접근 가능하게 만든다.
	
	\subsection{시나리오 B의 물리적 해석}
	
	만약 \(\LUV \approx 9\) TeV라면, 조정 네트워크는 플랑크 규모가 아닌 전약 규모에서 "결정화"된다. 형상 인자 \(F(q^2) = \exp[-(q^2/\LUV^2)^{2/3}]\)는 수 TeV 이상의 가상 운동량을 억제하기 시작한다. 이것은 즉각적인 관측 가능한 결과를 가진다:
	
	\begin{enumerate}[label=(\arabic*)]
		\item \textbf{LHC에서의 드렐-얀 및 다이제트 단면적}은 표준 모형 예측과 비교하여 높은 불변 질량에서 사건의 부족을 보여야 한다. 억제는 \(m_{jj} \gtrsim 5\) TeV에서 현저해진다.
		\item \textbf{힉스 쌍 생성} 및 벡터 보손 산란은 형상 인자가 힉스 자기 결합을 수정함에 따라 비정상적인 에너지 의존성을 보여야 한다.
		\item \textbf{전약 정밀 관측량의 편차}가 예상될 수 있지만, LEP 측정이 \(\LUV\)보다 훨씬 낮은 규모를 탐색하기 때문에 아마도 작을 것이다.
	\end{enumerate}
	
	결정적으로, 이 시나리오에서 최소 길이는 \(\Lmin = 2\pi\hbar c/\LUV \approx 1.38\times 10^{-19}\) m이며, 이는 플랑크 길이 (\(\sim 10^{-35}\) m)보다 훨씬 크다. 이것은 시공간의 입자성이 양자 중력 효과가 아니라 전약 규모 현상이며, 아마도 힉스 메커니즘 자체와 관련되어 있음을 의미한다.
	
	\subsection{전약 규모와 우주 상수 사이의 프랙탈 다리}
	\label{subsec:bridge}
	
	만약 시나리오 B가 올바르고 기본 조정 규모가 \(\LUV \approx 8.96\) TeV라면, 우주론에 대한 깊은 연결이 나타난다. YUCT에서 진공은 단일 연속 매체가 아니라 조정 층의 프랙탈 계층이다. 규모 \(\LUV\)는 표준 모형 장이 조정 형상 인자를 경험하는 층을 특징짓는다. 그러나 더 깊은(적외선) 층이 존재하며, 이는 점점 더 낮은 에너지 규모에 해당한다. 우주 상수, 즉 관측된 암흑 에너지 밀도는 가장 깊은 적외선 층에 의해 생성되며, 그 규모 \(\Lambda_{\mathrm{IR}}\)는 모든 양자 장의 영점 에너지를 전체 프랙탈 계층에 걸쳐 적분할 때 관측된 값 \(\rho_{\Lambda} \approx 2.51 \times 10^{-47}\) GeV\(^4\)을 재현해야 한다는 요구 조건에 의해 설정된다(완전한 유도는 부록~M 참조).
	
	자외선 규모 \(\LUV\)와 적외선 규모 \(\Lambda_{\mathrm{IR}}\) 사이의 관계는 YUCT의 기본 스케일링 양자에 의해 지배된다:
	\begin{equation}
		q = \left(\frac{3}{2}\right)^{1/3} \approx 1.1447.
	\end{equation}
	두 규모는 정수 \(N\)개의 조정 층에 의해 분리된다:
	\begin{equation}
		\LUV = \Lambda_{\mathrm{IR}} \cdot q^{N}.
		\label{eq:bridge}
	\end{equation}
	
	관측된 우주 상수로부터 \(\Lambda_{\mathrm{IR}} \approx 0.0032\) eV가 추출된다(부록~M, 식 (M.47) 참조). \(\LUV \approx 8.96 \times 10^{12}\) eV를 사용하면,
	\begin{equation}
		N = \frac{\ln(\LUV / \Lambda_{\mathrm{IR}})}{\ln q}
		= \frac{\ln(8.96 \times 10^{12} / 0.0032)}{\frac{1}{3}\ln(3/2)}
		\approx \frac{35.564}{0.135155} \approx 263.1.
		\label{eq:N_layers}
	\end{equation}
	을 얻는다.
	
	가장 가까운 반정수로 반올림하면(YUCT의 대수적 루프가 페르미온 층에 대해 요구하는 대로), 기본 구조 수
	\begin{equation}
		\boxed{N = 263.0}.
	\end{equation}
	를 얻는다.
	
	이것은 놀라운 결과이다. 전약 규모(힉스 보손의 질량)와 우주 상수는 자유 매개변수 없이 조정 계층의 정확히 \(263\)개의 이산적 단계에 의해 연결된다. 우주 상수의 관측된 작음은 미스터리가 아니다; 그것은 단순히 입자 물리학과 우주론을 분리하는 조정 네트워크의 엄청난 깊이를 반영할 뿐이다.
	
	\textbf{해석:} 계층의 각 층은 "규모 변환기" 역할을 한다: 주어진 층의 영점 에너지는 다음 더 깊은 층의 조정 형상 인자에 의해 억제된다. 총 진공 에너지는 모든 층의 기여의 합이며, 각 층은 이전 층보다 \(\beta = 2/3\)배 작다. 기하 급수는 유한한 값으로 수렴하며, 이는 정확히 관측된 암흑 에너지 밀도이다(자세한 합산은 부록~M 참조). 정수 \(N \approx 263\)은 전약 규모와 중성미자 질량 규모(가장 깊은 물리적으로 접근 가능한 층) 사이에 들어맞는 층의 수이다. 이 숫자가 반정수로 나타나는 것은 YUCT 프레임워크의 비자명한 일관성 검사이며, 시나리오 B를 더욱 지지한다.
	
	\subsection{두 시나리오의 비교}
	
	\begin{table}[H]
		\centering
		\caption{시나리오 A와 시나리오 B의 비교.}
		\label{tab:scenarios}
		\small
		\begin{tabular}{@{}lcc@{}}
			\toprule
			\textbf{속성} & \textbf{시나리오 A} & \textbf{시나리오 B} \\
			\midrule
			자외선 규모 \(\LUV\) & \(3\pi\Mpl \approx 10^{20}\) GeV & \(\approx 8.96\) TeV \\
			최소 길이 \(\Lmin\) & \(\frac{2}{3}\Lpl \approx 10^{-35}\) m & \(\approx 10^{-19}\) m \\
			라그랑지안의 맨 질량 & 0이 아님 (미세 조정) & 0 (규모 불변성) \\
			복사 보정 \(\delta m_H\) & \(10^{18}\) GeV & 125 GeV \\
			실험적 접근성 & 플랑크 규모 (접근 불가) & LHC 에너지 (접근 가능) \\
			미세 조정 필요 & 예 (\(10^{16}\) 분의 1) & 아니오 \\
			\bottomrule
		\end{tabular}
	\end{table}
	
	% ============================================================
	\section{시나리오 B의 반증 가능한 예측}
	\label{sec:predictions}
	
	만약 시나리오 B가 올바르다면, 다음 정량적 예측을 LHC와 미래 충돌기에서 테스트할 수 있다:
	
	\begin{enumerate}[label=(\arabic*)]
		\item \textbf{드렐-얀 단면적 억제.} 측정된 드렐-얀 단면적과 표준 모형 단면적의 비율은 딜렙톤 불변 질량 \(m_{\ell\ell}\)의 함수로서
		\[
		R(m_{\ell\ell}) \approx \exp\!\Bigl[-2\Bigl(\frac{m_{\ell\ell}^2}{\LUV^2}\Bigr)^{\!2/3}\,\Bigr],
		\]
		을 따라야 하며, 여기서 \(\LUV \approx 8.96\) TeV이다. \(m_{\ell\ell} \approx 5.4\) TeV에서 10\%의 억제가 예측되며, \(m_{\ell\ell} \approx 10\) TeV에서 \(\sim 2\)배로 증가한다.
		
		\item \textbf{힉스 쌍 생성.} 형상 인자는 높은 운동량 전달에서 힉스 자기 결합을 수정한다. 이-힉스 불변 질량 분포는 \(m_{HH} \gtrsim 3\) TeV에서 표준 모형으로부터 벗어나야 한다.
		
		\item \textbf{페르미온 질량.} 만약 모든 페르미온 질량도 방사적으로 생성된다면, 그들의 유카와 결합은 복사 보정이 관측된 질량 스펙트럼을 재현하도록 선택되어야 한다. 이것은 페르미온 조정 깊이 \(L_f\) 및 공명 인자 \(\xi_f\)를 복사 보정과 관련시키는 자기 일관성 조건을 부과한다. 자세한 분석은 향후 연구로 미룬다.
	\end{enumerate}
	
	\subsection{미세 구조 상수에 대한 주석}
	
	위상 불변량 \(N_{\mathrm{gauge}} = 12\)(부록~G)로부터 \(\alpha^{-1} \approx 137.036\)의 유도는 \(\LUV\)의 값과 무관하다. 게이지 결합은 대수적으로 흐른다; \(\LUV\)에서의 초기 값은 \(F_{18}\)의 위상에 의해 고정된다. 시나리오 A에서 \(\LUV \sim \Mpl\)이며, 플랑크 규모에서 전약 규모까지의 대수적 흐름은 \(\alpha^{-1} \approx 137.036\)을 제공한다. 시나리오 B에서 \(\LUV \sim 9\) TeV이며, 흐름은 훨씬 더 짧은 구간에 걸쳐 일어난다. 그러면 \(\alpha^{-1}\)에 대한 위상적 예측은 TeV 규모에 적용되며, 저에너지 값은 작은 대수적 진화에 의해 얻어질 것이다. 이것은 \(\alpha^{-1}\)의 수치적 예측을 수정할 것이며 CODATA 값과 비교되어야 한다. 예비 추정은 이동이 위상적 예측의 현재 불확실성 내에 있음을 시사하지만, 정확한 계산이 필요하며 부록 G의 미래 개정판에서 제시될 것이다.
	
	% ============================================================
	\section{논의: YUCT는 모든 것의 이론인가, 아니면 전약 규모의 이론인가?}
	\label{sec:discussion}
	
	원래 YUCT 프레임워크는 단일 조정 계층을 통해 플랑크 질량에서 전자 질량까지의 모든 물리적 규모를 설명하려고 했다. 수정된 스칼라 질량 계산은 이 야망 내의 긴장을 드러낸다: 만약 조정 네트워크가 플랑크 규모에서 작동한다면, 복사 보정은 거대하며 미세 조정된 맨 질량에 의해 상쇄되어야 한다.
	
	시나리오 B는 YUCT 질량 사다리를 재해석함으로써 이 긴장을 해결한다. 깊이 \(L\)을 통해 직접 질량을 생성하는 대신, 사다리는 \emph{유카와 결합}을 결정하며, 이것이 다시 \(\LUV \sim 9\) TeV의 규모에서 질량을 방사적으로 생성한다. 그러면 플랑크 규모는 조정 네트워크가 \emph{형성되기 시작하는} 규모가 되지만, 물리적 입자성 – 형상 인자가 중요해지는 규모 – 은 훨씬 더 낮다.
	
	이 재해석은 미세 조정 문제를 제거하면서 YUCT의 대수적 아름다움을 보존한다. 또한 YUCT를 사변적 양자 중력 이론에서 현재 및 미래 충돌기에서 테스트 가능한 반증 가능한 전약 물리학 모델로 변환한다.
	
	% ============================================================
	\section{결론}
	\label{sec:conclusion}
	
	우리는 YUCT 내의 스칼라 질량 보정에서 계산 오류를 수정하여 정확한 해석 값 \(C_0 = 3\sqrt{\pi}/(8\sqrt{2}) \approx 0.46999\)을 얻었다. 플랑크 규모 동일시 \(\LUV = 3\pi\Mpl\)를 사용하면, 힉스 질량에 대한 복사 보정은 \(\delta m_H \approx 1.57\times 10^{18}\) GeV이며, 이는 맨 질량에 대한 미세 조정된 상쇄를 요구한다.
	
	대안으로, 우리는 시나리오 B를 제안했다. 이 시나리오에서 플랑크 규모의 맨 라그랑지안은 규모 불변이며 모든 질량이 방사적으로 생성된다. 자기 일관성은 기본 조정 규모가 \(\LUV \approx 8.96\) TeV – LHC에서 접근 가능한 값 – 일 것을 요구한다. 이 시나리오는 고에너지 산란 과정에 대한 구체적이고 반증 가능한 예측을 제공한다.
	
	시나리오 A(미학적으로 만족스럽지는 않지만 YUCT의 플랑크 규모 기원과 일관됨)와 시나리오 B(예측 가능하고 테스트 가능하지만 YUCT 질량 사다리의 재해석을 요구함) 사이의 선택은 이론적 근거만으로는 이루어질 수 없다. 그것은 실험에 의해 결정되어야 한다. 우리는 LHC 협력단이 Eq.~(\ref{eq:F})에서 \(\LUV \sim 9\) TeV로 예측된 형상 인자 억제의 증거를 찾기 위해 고질량 드렐-얀 및 다이제트 데이터를 조사할 것을 촉구한다. 귀무 결과는 시나리오 B를 배제할 것이며, 시나리오 A는 논리적으로 가능하지만 계층 문제에 대한 해결책으로서 YUCT의 입지를 약화시킬 것이다.
	
	\section*{감사의 글}
	저자는 비판적 논의를 위해 YUCT 세미나 참가자들에게 감사드린다.
	
	\section*{데이터 가용성}
	모든 계산은 자체 완결적이다; 수치 평가에 사용된 Python 코드는 YPSDC 저장소 \url{https://github.com/Alexey-Yakushev-YUCT/YPSDC}에서 이용 가능하다.
	
	\begin{thebibliography}{99}
		\bibitem{AppendixL} A.\,V.\,Yakushev, ``부록 L: 프랙탈 조정 오차 스케일링,'' in \emph{YUCT}, Zenodo (2026).
		\bibitem{AppendixV} A.\,V.\,Yakushev, ``부록 V: 조정 과정의 엔트로피로서의 시간,'' in \emph{YUCT}, Zenodo (2026).
		\bibitem{AppendixG} A.\,V.\,Yakushev, ``부록 G: 조정 네트워크의 기하학적 장력으로서의 중력과 YUCT의 순환 우주론,'' in \emph{YUCT}, Zenodo (2026).
		\bibitem{AppendixM} A.\,V.\,Yakushev, ``부록 M: YUCT 우주론 – 조정 상전이로서의 인플레이션,'' in \emph{YUCT}, Zenodo (2026).
		\bibitem{AppendixLambda} A.\,V.\,Yakushev, ``부록 \(\Lambda\): YUCT 내 계층 문제와 우주 상수 문제의 해결,'' in \emph{YUCT}, Zenodo (2026).
		\bibitem{AppendixJ} A.\,V.\,Yakushev, ``부록 J: YUCT 위상적 오프셋으로부터 표준 모형 맛 매개변수의 완전한 유도,'' in \emph{YUCT}, Zenodo (2026).
		\bibitem{AppendixFerra} A.\,V.\,Yakushev, ``부록 Ferra1.2: 질서상 및 무질서상을 위한 보편적 조정 상수,'' in \emph{YUCT}, Zenodo (2026).
	\end{thebibliography}
	
\end{document}